\section{Discussion and Limitations}

The experimental results validate the core hypothesis that the composition of multiple reasoning patterns is superior to monolithic agentic loops for autonomous security operations.

\subsection{Insights into Pattern Synergy}
Our results highlight a critical synergy between planning and reflection. While ReWOO provides the efficient "skeleton" for the attack chain, the Reflexion engine serves as a vital reliability guardrail. The ablation study shows that disabling reflection reduces the success rate by 12.8\%, as the agent fails to recover from unforeseen tool errors or infrastructure nuances. Furthermore, the strategic lookahead provided by Tree of Thoughts (ToT) is essential for navigating the "Stealth vs. Speed" tradeoffs inherent in modern red teaming.

\subsection{Efficiency Scaling Laws}
We observe that the efficiency gains of CMatrix scale linearly with the depth of the attack chain. While ReAct agents must re-invoke the LLM after every tool observation, the CMatrix ReWOO planner batches observations, leading to a 2.42$\times$ reduction in token usage. This scaling behavior makes CMatrix uniquely suited for long-horizon enterprise assessments that would otherwise be cost-prohibitive.

\subsection{Governance as an Enabler}
The HITL safety gate demonstrates that formal governance can be an enabler rather than a bottleneck. By using a risk-aware interceptor, CMatrix maintains high throughput for non-disruptive tasks (92\% autonomy) while ensuring that 100\% of high-risk actions are human-validated. This architectural choice provides the necessary trust for deploying autonomous agents in production environments.

\subsection{Limitations and Future Work}
Despite these advancements, CMatrix faces limitations in extremely deep reasoning chains (15+ steps) where state context pollution can lead to increased hallucination. Future work will investigate hierarchical context pruning and the use of dynamic risk policies that adapt to the sensitivity of the target environment.
